\section{Requirements Review Procedure and Case Evidence}
\label{app:requirements-review}
Algorithm~\ref{alg:requirements-fidelity} details Figure~\ref{fig:requirements-review}. Campaign record $\mathcal{C}$ retains the English inputs, requirement table, candidates, costs, evidence, and decisions. Task $T$ names the reviewed scope, checks, and limits; $R$ identifies selected rule revisions. The table maps each source obligation to its interpretation, representation, assumptions, evidence, and review status.

\begin{figure}[!htbp]
\centering
% ARc-inspired staged workflow. The dashed path is proposed.
\begingroup
\hyphenpenalty=10000
\exhyphenpenalty=10000
\begin{tikzpicture}[
  x=1cm,y=1cm,font=\sffamily\fontsize{9}{10.6}\selectfont,
  card/.style={rounded corners=2pt,line width=.65pt,align=center,
    text width=2.72cm,inner xsep=3.5pt,inner ysep=3pt,minimum height=1.68cm},
  support/.style={card,minimum height=1.03cm},
  flow/.style={-{Stealth[length=1.7mm,width=1.2mm]},line width=.8pt,draw=black!85},
  advisory/.style={dashed,draw=scpBlue}
]
\node[card,scpRecord] (english) at (0,0)
  {\textbf{Natural language}\\[2pt]Reviewed\\requirements};
\node[card,scpProposal] (formal) at (3.54,0)
  {\textbf{Formalization}\\[2pt]SysML v2 /\\GUMBO};
\node[card,scpProposal] (learn) at (7.08,0)
  {\textbf{Supervised\\self-evolution}\\[2pt]Rule reuse\\and revision};
\node[card,scpCheck] (hamr) at (10.62,0)
  {\textbf{HAMR toolchain}\\[2pt]Generation and\\formal checks};
\draw[flow] (english) -- (formal);
\draw[flow] (formal) -- (learn);
\draw[flow] (learn) -- (hamr);

\node[support,advisory,fill=scpBlue!4,text width=5.82cm] (predict) at (1.63,-2.34)
  {\textbf{Proposed predictor}\\[2pt]Semantic agreement and next-attempt cost};
\node[support,scpRecord,text width=5.38cm] (review) at (9.25,-2.34)
  {\textbf{Human review and archive}\\[2pt]Meaning, scope, decisions,\\and check results};
\draw[flow,advisory] (predict.north -| formal.south) -- (formal.south);
\draw[flow] (hamr.south) -- (review.north -| hamr.south);
\draw[flow,draw=scpGray] (review.north -| learn.south) -- (learn.south);
\node[align=center,font=\sffamily\fontsize{8}{9.3}\selectfont] at (8.85,-1.30)
  {Diagnostics and\\reviewed outcomes};
\end{tikzpicture}
\endgroup

\caption{From reviewed natural-language requirements to checked artifacts, inspired by ARc's formalization and feedback loop~\cite[Fig.~1, p.~4]{bayless2025arc}. Diagnostics and reviewed outcomes inform repairs and rule updates. The dashed predictor would advise which candidate to check next. Actual allowance and recorded approvals govern attempts.}
\Description{Four stages link natural-language requirements, SysML v2 and GUMBO formalization, supervised rule use and revision, and HAMR generation and formal checks. A dashed proposed predictor estimates semantic agreement and next-attempt cost. Check results flow to human review and the archive. Diagnostics and reviewed outcomes feed supervised self-evolution.}
\label{fig:requirements-review}
\end{figure}

% Intended method composition; Figure requirements-review gives the main account.
% Historical case audits and the F01 recording limitation are stated in the appendix.
\begin{algorithm}[!htbp]
\caption{Review requirements fidelity during development}
\label{alg:requirements-fidelity}
\KwIn{Campaign record $\mathcal{C}$; task $T$; selected rule revisions $R$.}
\KwOut{Updated $\mathcal{C}$ and scoped task outcome $q$.}
Retain the original English, input versions, approved scope, and rule revisions\nllabel{ln:rf-retain}\;
Read the source line by line; record obligations, interpretations, representations, and assumptions in the requirement table\;
Assess meaning and unresolved questions; record the responsible reviewer or agent\;
\lIf{$T$ lacks current human approval}{\Return \scpStatus{pending}}
\While{authority, deadline, and allowance after consumption and reservations permit an attempt}{
  \lIf{usage is unresolved}{\Return \scpStatus{pending}}
  \lIf{reviewed inputs or dependencies have changed}{\Return \scpStatus{stale}}
  Reserve allowance; request a proposal for $T$ using $R$\;
  Record usage, proposal, and detailed call outcome\;
  \lIf{response is invalid or check time is unavailable}{\Return \scpStatus{pending}}
  Run HAMR and the required checks; retain results, failures, obligation links, and content fingerprints\nllabel{ln:rf-check}\;
  Update the requirement table; review interpretation, assumptions, and unresolved obligations\;
  Apply current human decisions with Algorithm~\ref{alg:attestation}; request any required new review; obtain $q$\;
  \If{$q$ is accepted}{
    Mark the included task scope complete; record the owner's release decision when supplied, with exclusions and remaining obligations\;
    \Return $q$\;
  }
  \lIf{$q$ is stale or rejected}{\Return $q$}
  \If{the next action changes meaning, scope, assumptions, or limits}{
    Retain the proposal; \Return \scpStatus{pending} for human review\nllabel{ln:rf-amend}\;
  }
  \lIf{rule revision follows a current failed \texttt{NOT\_VALIDATED} use of the same rule and learning plan}{\Return \texttt{NOT\_VALIDATED}}
  \lIf{the next action lacks repair authority}{\Return $q$}
  Retain earlier charges and failures; refine within the approved scope and retry\;
}
\Return \scpStatus{pending} for further allowance or instruction\;
\end{algorithm}


This procedure composes development, accounting, and review services~\cite{scpCheckpoint}. The reported Producer/Consumer and SNG cases used saved source and proof audits. Their review criteria followed the start of generation; SNG's final six-relation audit followed runtime. Appendix~\ref{app:coverage-details} gives the actual mappings and service configuration.

Appendix~\ref{app:semantic-budget} describes the proposed use of predictions to select attempts. Historical engineering results retain their status; later owner release decisions record their dates and scopes.
